\section{Introduction}
\label{sec:introduction}

If domestic money creation mechanically drove inflation, why did Central Bank base money expand by over 100\% in 1971 while monthly inflation remained quiescent at 1.5--2.0\% and annual inflation fell from 35\% to 22\%? For over half a century, the macroeconomic crisis of Chile's Unidad Popular (UP, 1970--1973) under Salvador Allende has been treated by orthodox economics as a cautionary tale of populist fiscal and monetary mismanagement \citep{LarrainMeller1991,Edwards2024}. Under closed-economy assumptions and monetarist quantity theory, this literature attributes the 1972--1973 price explosion directly to discretionary central bank emission. Yet the 1971 record presents an immediate empirical puzzle for that premise. The non-inflationary absorption of rapid monetary expansion in 1971 rested on balance-sheet solvency and real capacity: the Central Bank of Chile (BCCh) held nearly \$400~million in inherited foreign exchange reserves, and industrial firms mobilized idle capacity, with structural capacity utilization reaching 98.4\%. Significant inflation emerged only in mid-1972, after international reserves were exhausted and a targeted external credit cutoff severed revolving commercial trade lines. Domestic money growth functioned primarily as an endogenous accommodation valve through which the monetary authority mediated an acute balance-of-payments squeeze amidst intense distributive conflict.

This paper argues that the transmission of domestic liquidity was governed by the central bank's command over foreign exchange. Grounding dependency theory in open-economy balance sheets, I formalize the Central Bank's position within the international currency hierarchy \citep{Vasudevan2009,Basu2022}. The central bank balance sheet in a dependent economy embodies a fundamental structural asymmetry: on the liability side, domestic base money ($H_t$) consists of sovereign credit tokens issued under national authority; on the asset side, international reserves ($F_t$) represent claims on world money that the peripheral state cannot emit. To evaluate how reserve backing conditions monetary transmission, I define the Solvency Ratio as convertible foreign reserves over base money ($S_t \equiv F_t / H_t$) and track its dynamic evolution via the monthly Solvency Growth Gap ($\Delta s_{t-1} \equiv g^F_{t-1} - g^H_{t-1}$). Threshold econometric estimates identify a structural transition into an adverse reserve-coverage regime when base money growth outpaces foreign reserve growth by more than 4.6 percentage points ($\Delta s_{t-1} \le -4.615\%$).

To evaluate these competing hypotheses across distinct analytical horizons, the empirical design deploys a two-tier time-series architecture. The annual tier ($1920$--$2010$, $T=91$) combines functional income shares from \citet{Astorga2023}, balance-of-payments accounts from \citet{NievasPiketty2025}, and net capital stock estimates from the Perpetual Inventory (GPIM) dataset developed in \citet{Polanco2026}. At annual frequency, the Austrian malinvestment hypothesis \citep{EspinosaVejar2025} fails to find empirical support: gross fixed investment in machinery and equipment contracted ($-2.3\%$) during 1971 while capacity utilization surged to an all-time peak of $98.4\%$, consistent with idle capacity reactivation rather than intertemporal capital misallocation. To isolate dynamic transmission at high frequency, the monthly tier ($1960$--$1980$, $N=250$ continuous months) integrates Central Bank statistical database series on monetary base, gold, copper, and international prices, alongside activity indicators compiled by \citet{DiazBahamonde2023} and monthly IMF reserve data. These series document that Prebisch's purchasing power capacity to import ($M^{\text{cap}}_t$) collapsed following the 1971 international monetary shock, driving the Real Structural Imbalance Ratio ($\Theta_t \equiv M_t / M^{\text{cap}}_t \times 100$) to a peak of $283.4$ in May 1973.

The econometric results establish three central empirical findings. First, directional precedence tests in stationary vector autoregressions demonstrate that reverse defensive accommodation dominates forward transmission: consumer price inflation predicts subsequent base-money growth across all tested horizons ($k=1,\dots,4$) with $F$-statistics approximately four times larger than forward money-to-price predictability. In the adverse reserve-coverage regime of the Threshold VAR, generalized impulse response functions (GIRFs) show that reverse accommodation ($\pi_m \to g^H$) produces an immediate response peaking at 1.72 percentage points at month 1 and remaining statistically significant for 10 consecutive months, delivering a cumulative 12-month expansion of 6.54 percentage points (point estimate; 7.00 pp bootstrap median). By comparison, forward transmission ($g^H \to \pi_m$) yields a cumulative 12-month price increase of 3.55 percentage points (point estimate; 3.90 pp bootstrap median). Second, the external reserve buffer governs international price pass-through: in the non-adverse regime, international price shocks ($g^{\text{gold}} \to \pi_m$) have no statistically significant effect on domestic inflation over a 12-month horizon; once reserve coverage deteriorates ($\Delta s_{t-1} \le -4.615\%$), the same shock generates an unbroken 10-month inflationary wave yielding 6.54 percentage points of cumulative inflation (point estimate; 6.93 pp bootstrap median). Dynamic multiplier difference tests ($\Delta(h) \equiv \text{GIRF}^{\text{Insolv}}(h) - \text{GIRF}^{\text{Normal}}(h)$) confirm statistically significant cross-regime divergence across 10 consecutive months ($h=3,\dots,12$). Third, both forward and reverse transmissions intensify under reserve stress: without foreign reserves to anchor the currency and finance imported inputs, domestic cost-push pressures and credit expansions operate unbuffered.

These findings allow an empirical discrimination between rival research programmes in Chilean economic thought. Monetarist interpretations \citep{Edwards2024} treat high-powered money as an exogenous policy lever and assume direct, regime-invariant pass-through to prices under full-capacity market clearing. Austrian accounts \citep{EspinosaVejar2025} interpret the crisis as credit-induced capital malinvestment. In contrast, dependency theory \citep{Prebisch1950,Marini1973,Kvangraven2021} and modern International Political Economy \citep{Vasudevan2009,Basu2022} locate the primary constraint in external financial subordination and unequal access to world money. Evaluated against the sign-discrimination framework in Section~\ref{sec:literature_review}, orthodox models fail to account for the 1971 non-inflationary expansion, the strong precedence of prices over money, and the sharp regime asymmetry in external shock transmission. The open-economy balance-sheet model explains these anomalies as direct consequences of peripheral reserve dynamics.

The Chilean transition unfolded during the global structural crisis of the early 1970s. President Richard Nixon's unilateral suspension of dollar-gold convertibility on August 15, 1971, dismantled the Bretton Woods par-value system and transmitted stagflationary shocks across the periphery \citep{Vernengo2021,Zoeller2019}. This global dislocation interacted with a targeted bilateral credit squeeze: foreign commercial banks canceled over \$220~million in revolving trade credit lines, while the U.S. Eximbank and multilateral lenders halted disbursements \citep{DeVylder1976,PetrasMorley1975}. As documented in the Central Bank's \textit{Memorias Anuales} and \textit{Boletines Mensuales}, the monetary authority exercised constrained institutional agency within this hostile external environment. Facing enterprise illiquidity, spare parts shortages, and distribution bottlenecks, the Central Bank expanded the discount window defensively to finance essential food imports, maintain firm solvency, and sustain worker payrolls. Rather than initiating an autonomous monetary experiment, the Central Bank accommodated structural cost shocks to prevent production collapse under an acute foreign exchange blockade.

This paper extends the political economy of peripheral accumulation and external vulnerability developed in \citet{Polanco2026}. While \citet{Polanco2026} established the annual balance-of-payments boundary throttling mechanization under unbalanced growth, this paper investigates the dynamic balance-sheet mechanisms through which that external constraint operates at monthly frequency. The remainder of this paper is structured as follows. Section~\ref{sec:literature_review} reviews the competing analytical traditions and specifies the sign-discrimination framework. Section~\ref{sec:macro_framework} develops the theoretical model of world money, financial subordination, and central-bank balance-sheet mechanics. Section~\ref{sec:data_sources} details the data compilation and variable construction. Section~\ref{sec:historical_empirical_results} presents the empirical results across the linear VAR Granger tournament and the non-linear Threshold VAR system. Section~\ref{sec:discussion_conclusion} discusses the political economy implications, addresses methodological limits, and draws strategic conclusions. Methodological derivations, archival codebooks, and complete impulse response trajectories are provided in Appendices~\ref{app:bop_levr}--\ref{app:linear_var_diagnostics}.
